\documentclass[11pt]{article}

\usepackage{amsmath,amssymb}

\title{Problem Statement and Normalization Conventions}
\author{}
\date{}

\begin{document}
\maketitle

Fix a positive integer $n$.  The set of admissible Seidel matrices of
order $n$ is
\[
  \mathcal{S}_n
  :=
  \left\{
    A=(a_{ij})\in\mathbb{R}^{n\times n}:
    A^{\mathsf T}=A,\quad
    a_{ii}=0,\quad
    a_{ij}\in\{-1,1\}\ \text{for }i\ne j
  \right\}.
\]
The sign cube is
\[
  {\{-1,1\}}^{n}
  =
  \left\{x={(x_1,\ldots,x_n)}^{\mathsf T}:x_i\in\{-1,1\}\right\}.
\]

For $A\in\mathcal{S}_n$ and $x\in{\{-1,1\}}^{n}$, define the pair-sum
quadratic by
\[
  Q_A(x):=\sum_{1\le i<j\le n}a_{ij}x_i x_j.
\]
Because $A$ is symmetric and has zero diagonal, the full quadratic form is
\[
  x^{\mathsf T}Ax
  =\sum_{i=1}^n\sum_{j=1}^n a_{ij}x_i x_j
  =2\sum_{1\le i<j\le n}a_{ij}x_i x_j
  =2Q_A(x),
\]
or, equivalently, $Q_A(x)=\tfrac12 x^{\mathsf T}Ax$.

For a fixed admissible matrix, set
\[
  M(A):=\max_{x\in{\{-1,1\}}^{n}}\left|x^{\mathsf T}Ax\right|.
\]
Thus the two quadratic normalizations satisfy
\[
  M(A)
  =2\max_{x\in{\{-1,1\}}^{n}}|Q_A(x)|.
\]
The extremal quantity is
\[
  F(n)
  :=\min_{A\in\mathcal{S}_n}
       \max_{x\in{\{-1,1\}}^{n}}|Q_A(x)|
  =\frac12\min_{A\in\mathcal{S}_n}M(A),
\]
and its normalized form is
\[
  L_n:=\frac{F(n)}{n^{3/2}}.
\]

The open question is: does the limit
\[
  \lim_{n\to\infty}L_n
\]
exist, and if it exists, what is its value?

\end{document}
